% ============================================================
% NUMERICAL ACCURACY ASESSMENT
% ============================================================

\section*{Numerical Accuracy Assessment}

% ----------------------------------------------------------
% Classical Thickness Distributions
% ----------------------------------------------------------

\subsection*{Scope of the Error Analysis}

The present tool does not involve CFD or flow-field simulation.
The geometry is generated from analytical NACA formulations and, in some cases,
is reconstructed from tabulated coordinate data using polynomial interpolation.
Therefore, the primary sources of numerical error arise from:

\begin{enumerate}
    \item Discretization error due to finite chordwise sampling.
    \item Interpolation error from polynomial reconstruction.
    \item Tabulated data approximation and rounding.
\end{enumerate}

Since the analytical NACA equations are available, 
they are treated as the reference solution for validation purposes.

% ----- Discretization Error -----

\subsubsection*{Discretization Error}

Let $y_{exact}(x)$ denote the analytical airfoil surface and
$y_N(x)$ the discretized representation using $N$ chordwise points.

The discretization error is evaluated using:
\[
E_{\infty}(N) = \max_{i} | y_N(x_i) - y_{exact}(x_i) |
\]

and the root-mean-square (RMS) error:
\[
E_{RMS}(N) =
\sqrt{
\frac{1}{N}
\sum_{i=1}^{N}
\left(
y_N(x_i) - y_{exact}(x_i)
\right)^2
}
\]

A grid refinement study is performed by increasing
$N = 50, 100, 200$.

For smooth analytical airfoil functions,
the discretization error is expected to decrease approximately with
$\mathcal{O}(N^{-2})$ for sufficiently smooth regions.

% ----- Effect of Node Distribution -----

\subsubsection*{Effect of Node Distribution}

Uniform spacing and cosine spacing are compared under identical $N$.

Since airfoil curvature is highest near the leading edge, 
uniform spacing under-resolves this region, 
leading to locally amplified errors.

Cosine spacing clusters nodes near $x=0$, 
thus significantly reducing leading-edge geometric error.
For typical NACA 4-digit profiles, 
the maximum geometric deviation near the leading edge is expected to decrease by one order of magnitude when cosine spacing is used instead of uniform spacing at moderate resolution ($N \approx 100$).

% ----- Interpolation Error -----

\subsubsection*{Interpolation Error}

When reconstructing geometry from tabulated data, 
Lagrange polynomial interpolation of degree $n$ is applied.

The interpolation error is theoretically given by:
\[
f(x) - P_n(x) =
\frac{f^{(n+1)}(\xi)}{(n+1)!}
\prod_{i=0}^{n}(x - x_i)
\]

for some $\xi$ in the interpolation interval.

Since airfoil thickness and camber distributions are smooth functions, 
the interpolation error remains small for moderate $n$.
However, high-degree interpolation under uniform spacing may produce oscillations (Runge phenomenon).

In practical implementation,
interpolation is restricted to a moderate polynomial degree or applied locally to avoid instability.

The expected interpolation error magnitude,
for $n \leq 10$ and smooth tabulated data,
is typically below $10^{-4}c$ in geometric coordinates.

% ----- Tabulated Data Approximation -----

\subsubsection*{Tabulated Data Approximation}

Tabulated airfoil data are usually rounded to
$4$--$5$ decimal places.
This introduces a baseline numerical uncertainty
on the order of $E_{table} \sim 10^{-5}c \text{ to } 10^{-4}c$, which represents the lower bound of achievable accuracy
when reconstructing geometry from such data.

% ----- Overall Expected Accuracy -----

\subsubsection*{Overall Expected Accuracy}

Combining discretization and interpolation effects,
the total geometric error may be estimated as:
\[
E_{total} \approx
E_{disc} + E_{interp} + E_{table}
\]

For typical configurations with $N \geq 100$ cosine-spaced nodes and a moderate interpolation degree, 
the overall geometric deviation is expected to remain within $E_{total} \leq 10^{-4}c$, 
corresponding to less than $0.01\%$ of the chord length.

% ----- Acceptable Error and Adjustment Criteria -----

\subsubsection*{Acceptable Error and Adjustment Criteria}

Since the present tool is intended for geometric visualization,
preliminary aerodynamic estimation,
and educational or conceptual analysis,
a geometric tolerance of $E_{tol} = 10^{-4}c$
is considered acceptable.

If the error exceeds this threshold,
favorable adjustments are: 
increase chordwise resolution, 
use cosine spacing instead of uniform spacing, 
reduce the interpolation degree or switch to spline interpolation, 
refine tabulated input precision.

% ----- Conclusion -----

\subsubsection*{Conclusion}

The dominant numerical error source in the present implementation
is discretization near the leading edge.
Cosine spacing effectively mitigates this issue.

Under typical settings, the geometric error remains sufficiently small
for engineering-level preliminary airfoil analysis.
Further refinement is only necessary if high-fidelity CFD or structural analysis is to be performed using the generated geometry.